\section*{Overview}

Li Chao tree is the line-container I use when a DP or query problem reduces to
\[
\min_j (m_j x + b_j)
\]
or \(\max\) instead, but the pleasant monotonicity assumptions behind the deque version of the convex hull trick are
missing. It gives a robust contest implementation for arbitrary line insertion order and arbitrary query order.

\section*{Problem-Driven Motivation}

A very common contest transition looks like this:
\[
dp[i] = cost(i) + \min_{j < i}\bigl(A_j \cdot X_i + B_j\bigr).
\]

The naive implementation tries every previous \(j\), so it is \(O(n^2)\). That dies immediately when \(n\) is
\(2 \cdot 10^5\).

If both slopes \(A_j\) and query points \(X_i\) are monotone, a deque-based convex hull trick is enough. But many real
problems break one or both assumptions:

\begin{itemize}
  \item the \(x\)-coordinates come in arbitrary order,
  \item the states that generate lines are not sorted by slope,
  \item queries and insertions are interleaved online.
\end{itemize}

That is the point where Li Chao tree becomes the safe fallback.

\section*{Recognition Pattern}

Strong signals for Li Chao tree are:

\begin{itemize}
  \item a recurrence or query of the form ``minimum / maximum over lines evaluated at one \(x\),''
  \item arbitrary query order, arbitrary line insertion order, or both,
  \item a value domain that can be bounded in advance or compressed offline,
  \item the monotone CHT looks almost right but one monotonicity assumption is false.
\end{itemize}

If the compared functions are not lines, or if two functions can intersect many times, Li Chao is probably the wrong
tool.

\section*{Derivation}

Suppose every candidate state contributes one line \(y = mx + b\), and every new state asks for the best value at one
query point \(x\). The remaining task is:

\begin{quote}
Maintain the lower envelope of inserted lines and query it at points.
\end{quote}

The Li Chao idea is to put a segment-tree structure on the \(x\)-domain. Each node owns an interval \([l,r]\) and stores
one currently best candidate line for that interval.

The crucial observation is that two lines intersect at most once. So if we compare two lines on an interval and one of
them is better at the midpoint, then the worse line can only still matter on \emph{one side} of that midpoint.

\includegraphics[width=\textwidth]{figures/query-partition.svg}

That leads to the insertion rule:

\begin{itemize}
  \item keep the better line at the midpoint in the current node,
  \item recurse with the other line only into the side where it could still win.
\end{itemize}

This is why Li Chao stays logarithmic: every insertion follows only one path downward after each comparison.

\section*{Worked Problem}

\paragraph{Problem.}
Given arbitrary integers \(x_0, x_1, \dots, x_{n-1}\), compute
\[
dp[0] = 0,\qquad
dp[i] = x_i^2 + C + \min_{0 \le j < i}\bigl(dp[j] + x_j^2 - 2x_i x_j\bigr).
\]

\paragraph{Why the naive solution fails.}
Trying all \(j < i\) for every \(i\) is \(O(n^2)\). For \(n = 2 \cdot 10^5\), that is far too slow.

\paragraph{How the formula becomes lines.}
For fixed \(j\), the part depending on \(x_i\) is
\[
(-2x_j)\cdot x_i + (dp[j] + x_j^2).
\]

So each previous state contributes a line with:

\begin{itemize}
  \item slope \(m_j = -2x_j\),
  \item intercept \(b_j = dp[j] + x_j^2\),
  \item query point \(x = x_i\).
\end{itemize}

Then
\[
dp[i] = x_i^2 + C + \min_j (m_j x_i + b_j).
\]

\paragraph{Why Li Chao, not monotone CHT.}
The \(x_i\) values may arrive in any order, so front-popping from a deque is illegal. The slopes
\(-2x_j\) are also unsorted if the \(x_j\) values are unsorted. That rules out the monotone hull but not Li Chao.

\paragraph{Final algorithm.}

\begin{itemize}
  \item bound the queried \(x_i\) values by \([\min x_i, \max x_i]\),
  \item insert the line for state \(0\),
  \item for \(i = 1 \ldots n-1\), query the best line at \(x_i\), compute \(dp[i]\), then insert the line generated by
  state \(i\).
\end{itemize}

\section*{Implementation Reasoning}

The invariant I want from the implementation is:

\begin{quote}
For every coordinate \(x\), at least one line on the root-to-leaf path for \(x\) is optimal among all inserted lines.
\end{quote}

That is why queries simply evaluate every stored line on that path and take the best.

Practical implementation details:

\begin{itemize}
  \item use \texttt{\_\_int128} for comparisons, because \(mx+b\) can overflow 64-bit during intermediate arithmetic,
  \item handle equal slopes carefully: for a minimum hull, keep only the smaller intercept,
  \item if the domain is sparse but all query points are known offline, compress the coordinates and build Li Chao over
  compressed positions.
\end{itemize}

The code keeps a dynamic-node tree because it is memory-friendly when the coordinate domain is large but only a small
part of it is touched.

\section*{Correctness Intuition}

Inside one node interval, the kept line is the one that wins at the midpoint. Since two lines intersect at most once,
the line that loses at the midpoint can only still become relevant on one side. Recursing only into that side preserves
all possible future optima.

Therefore, after every insertion, the optimal line for any queried \(x\) still survives somewhere on the path from the
root to the leaf containing \(x\).

\section*{Complexity Analysis}

If the domain spans \(X\) discrete coordinates:

\begin{itemize}
  \item insert one line: \(O(\log X)\),
  \item query one point: \(O(\log X)\),
  \item memory: \(O(\text{insertions} \cdot \log X)\) in the dynamic-node version.
\end{itemize}

\section*{Common Pitfalls}

\begin{itemize}
  \item Forgetting that Li Chao only supports point queries unless you implement the segment-restricted variant.
  \item Choosing a coordinate range that misses some query points.
  \item Using it on functions that are not single-intersection objects.
  \item Flipping the inequalities incorrectly when switching between minimum and maximum.
  \item Assuming 64-bit multiplication is safe during comparisons.
\end{itemize}

\section*{Variants and Failure Modes}

\begin{itemize}
  \item If slopes and query points are both monotone, deque CHT is simpler and faster.
  \item If a line is active only on a subsegment of the domain, use the segment-restricted Li Chao variant.
  \item If the domain is continuous and precision matters, use a floating version carefully.
  \item If the compared functions intersect multiple times, Li Chao is no longer justified by the midpoint argument.
\end{itemize}

\section*{Practice Problems}

\begin{itemize}
  \item DP with arbitrary-order quadratic transitions.
  \item Dynamic minimum-cost linear functions evaluated at points.
  \item Segment-restricted linear cost updates after offline event expansion.
\end{itemize}

\section*{References}

\begin{itemize}
  \item \href{https://cp-algorithms.com/geometry/convex_hull_trick.html}{cp-algorithms: Convex Hull Trick and Li Chao Tree}.
  \item \href{https://github.com/kth-competitive-programming/kactl}{KACTL}.
  \item \href{https://codeforces.com/catalog}{Codeforces Catalog}.
\end{itemize}
